  \documentclass[12pt,a4paper]{ctexart}
  % 公共排版文件（使用 XeLaTeX 编译）
\usepackage[a4paper,margin=2.35cm]{geometry}
\usepackage{booktabs,longtable,tabularx,array,multirow,enumitem}
\usepackage{graphicx,xcolor,amsmath,hyperref,lastpage,fancyhdr}
\usepackage{tikz}
\usetikzlibrary{arrows.meta,positioning,shapes.geometric,fit,calc}
\hypersetup{colorlinks=true,linkcolor=blue!60!black,urlcolor=blue!60!black,citecolor=blue!60!black,pdftitle={逸仙活动云课程报告}}
\setlist{itemsep=2pt,topsep=3pt,leftmargin=2em}
\setlength{\parindent}{2em}
\setlength{\headheight}{15pt}
\renewcommand{\arraystretch}{1.35}
\newcolumntype{Y}{>{\raggedright\arraybackslash}X}
\newcommand{\reportcover}[1]{%
  \begin{titlepage}
  \begin{center}
  \vspace*{3cm}{\zihao{1}\bfseries 中山大学\\[0.6cm]《软件工程》期末大作业\par}
  \vspace{2.5cm}{\zihao{2}\bfseries #1\par}
  \end{center}
  \vspace{2cm}
  \begin{center}
  \begin{tabular}{l}
  {\Large 项目名称：逸仙活动云——中山大学校园活动云平台}\\
  {\Large 小组成员：钟梓俊\quad 邱剑波\quad 张宸睿}\\
  {\Large 指导老师：王可泽}
  \end{tabular}
  \end{center}
  \vfill
  \begin{center}
  {\large 2026 年 7 月 16 日\par}
  \end{center}
  \end{titlepage}}

\pagestyle{fancy}\fancyhf{}\lhead{逸仙活动云}\rhead{软件工程期末大作业}\cfoot{第 \thepage\ 页，共 \pageref{LastPage} 页}
\setcounter{tocdepth}{2}

  \begin{document}
  \reportcover{需求分析报告}
  \pagenumbering{Roman}\tableofcontents\clearpage\pagenumbering{arabic}

  \section{项目概述}
  \subsection{项目背景}
  校园活动是高校学生学习生活的重要组成部分。然而，活动信息分散于学院网站、学生组织公众号、通知群、线下海报等多个渠道，发布时间不统一、格式不一致，学生需要反复在不同平台间查找。活动临近时，报名入口遗失、时间冲突、详情不完整等问题也时常出现。对活动组织者而言，活动发布、审核、报名统计和通知推送缺少统一平台支撑，重复劳动多，管理经验难以沉淀。

  本项目“逸仙活动云”面向校园活动信息的发现、发布、审核、参与和治理场景，建设一个 Web 端校园活动信息平台。系统前端采用 Vue 3、TypeScript、Vite、Pinia、Vue Router 和 Element Plus，后端采用 Flask、SQLAlchemy、JWT、蓝图与服务层结构，并提供 SQLite 本地运行路径以及 PostgreSQL、Redis、Celery 的容器化部署入口。

  \subsection{项目目标}
  项目目标概括为“让活动被看见、被管理、被可信地参与”。具体包括：
  \begin{enumerate}
  \item 为学生提供统一的活动浏览、检索、报名、收藏、日历和订阅入口，降低发现活动和管理个人安排的成本。
  \item 为活动发布者提供活动草稿、提交审核、审核反馈查看和活动维护能力，减少多渠道重复发布。
  \item 为管理员提供发布者申请审批、活动审核、数据源管理、抓取日志和审计日志查看能力，保证平台内容可信、操作可追溯。
  \end{enumerate}

  \subsection{项目范围}
  项目范围界定如下：
  \begin{enumerate}
  \item \textbf{功能范围：} 系统覆盖活动发现与检索、身份认证与角色授权、活动发布与审核、个人参与管理（报名、收藏、日历）、发布者申请管理、知识关联与订阅通知、数据源与内容质量管理、审计与安全治理八项功能域（详见第 3 节）。
  \item \textbf{系统范围：} 交付物包括 Vue 3 前端单页应用、Flask REST API 后端、本地 Mock 联调服务以及 Docker Compose 容器化部署拓扑。
  \item \textbf{约束条件：} 采用关系型数据库保证数据一致性；JWT + RBAC 实现服务端授权，前端不以按钮隐藏替代后端权限校验；外部邮件、搜索和 AI 服务为可选依赖，不可用时核心浏览、认证与审核流程须能降级运行；前端支持 360px 以上视口宽度。
  \item \textbf{排除项：} 移动原生应用、校级统一身份认证、在线支付、生产级监控告警平台、多语言国际化等不属于本次项目范围。
  \end{enumerate}

  \section{用户需求分析}
  \subsection{用户痛点}
  \begin{enumerate}
  \item 信息分散：活动分布在多个渠道，用户无法通过一个入口完整了解校园活动。
  \item 检索效率低：缺少统一分类、关键词搜索、状态筛选和排序能力，用户难以及时找到目标活动。
  \item 参与管理弱：报名、收藏、日历和提醒分散在不同工具中，学生容易忘记活动时间或重复报名。
  \item 发布成本高：组织者需要在不同平台重复编辑和发布活动，审核结果反馈也不集中。
  \item 内容可信度不足：来源、审核状态、重复内容和质量问题难以追踪，管理员缺少治理工具。
  \item 协作和演化困难：若缺少统一需求基线、接口约束和测试证据，后续迭代容易出现前后端字段不一致、权限遗漏和回归缺陷。
  \end{enumerate}

  \subsection{核心用户需求}
  核心用户需求以“用户价值”为中心描述：
  \begin{enumerate}
  \item 学生需要快速发现可信、完整、可筛选的校园活动信息。
  \item 学生需要把报名、收藏、日历和通知统一管理，避免遗漏重要活动。
  \item 发布者需要低成本提交活动，并清楚知道审核进度和处理结果。
  \item 管理员需要以角色权限控制平台内容质量，并能够追溯关键操作。
  \item 开发和运维人员需要通过清晰配置、自动化测试和部署文档保证项目可交付、可复现、可继续演化。
  \end{enumerate}

  \subsection{用户角色}
  根据调研，拟确定平台包含以下五种用户角色：\par
  \noindent\begin{tabularx}{\textwidth}{p{2.4cm}Y Y}
  \toprule 角色 & 角色说明 & 主要需求 \\\midrule
  访客 & 未登录用户，主要用于浏览公开活动。 & 查看活动列表、搜索活动、查看详情，并在需要报名或收藏时被引导登录。 \\
  普通用户 & 已注册登录的学生用户，系统默认角色为 \texttt{viewer}。 & 管理个人资料，报名/取消报名，收藏活动，加入个人日历，订阅感兴趣内容，接收通知。 \\
  发布者 & 经管理员审批后获得发布权限的用户，角色为 \texttt{publisher}。 & 创建活动草稿，编辑活动内容，提交审核，查看审核状态和驳回原因。 \\
  管理员 & 负责平台内容治理和权限管理，角色为 \texttt{admin}。 & 审核活动，审批发布者申请，维护数据源，查看抓取日志与审计日志，处理异常内容。 \\
  运维/开发人员 & 负责部署、配置、测试和故障排查。 & 使用环境变量配置系统，运行自动化测试，查看日志和健康检查结果，完成版本发布与回滚。 \\
  \bottomrule
  \end{tabularx}

  \subsection{用户使用场景}
  下表展示不同用户角色如何通过我们的平台解决上述存在的痛点：

  \begin{longtable}{p{2.2cm}p{3.2cm}p{6.8cm}}
  \toprule 场景 & 触发条件 & 期望结果 \\\midrule\endhead
  活动发现 & 学生希望了解近期讲座、比赛或社团活动。 & 用户可在首页或活动列表中按关键词、分类、状态和排序方式筛选，并进入详情页查看时间、地点、主办方和正文。 \\
  活动参与 & 用户决定参加某个活动。 & 登录后完成报名，可取消报名，可收藏活动，并可将活动加入个人日历或导出 ICS。 \\
  内容发布 & 发布者需要发布一项新活动。 & 发布者先保存草稿，补全必要信息后提交审核，管理员审核通过后活动对外展示。 \\
  审核治理 & 管理员发现待审核活动或发布者申请。 & 管理员查看详情，选择通过或驳回并填写原因，系统记录审计日志。 \\
  订阅通知 & 用户关注某类知识节点或关键词。 & 当新活动匹配订阅条件时生成站内通知，用户可查看并标记已读。 \\
  异常处理 & 网络失败、未登录、越权或数据为空。 & 前端展示明确的加载、空态、错误提示或登录引导，后端返回稳定 JSON 错误信息。 \\
  \bottomrule
  \end{longtable}



  \section{功能需求分析}
  \subsection{功能模块划分}
  \noindent\begin{tabularx}{\textwidth}{p{2.7cm}Y Y}
  \toprule 模块 & 主要功能 & 面向角色 \\\midrule
  活动发现模块 & 首页活动推荐、活动列表、关键词搜索、分类/状态筛选、活动详情。 & 访客、普通用户 \\
  身份认证模块 & 图形验证码登录、邮箱验证码注册、JWT 会话、用户资料、角色鉴权。 & 访客、普通用户、发布者、管理员 \\
  个人参与模块 & 报名、取消报名、收藏、取消收藏、个人日历、ICS 导出。 & 普通用户 \\
  活动发布模块 & 草稿创建、活动编辑、提交审核、查看审核结果。 & 发布者、管理员 \\
  审核管理模块 & 活动审核、批量审核、发布者申请审批、审核意见记录。 & 管理员 \\
  知识与通知模块 & 知识节点关联、订阅管理、通知生成、已读标记。 & 普通用户、发布者、管理员 \\
  数据治理模块 & 数据源维护、抓取日志、内容质量分、疑似重复标识、审计日志。 & 管理员、运维人员 \\
  异常与通用模块 & 加载态、空态、错误态、401/403 处理、接口统一错误响应。 & 全部角色 \\
  \bottomrule
  \end{tabularx}

  \subsection{核心功能需求}
  \begin{longtable}{p{1.2cm}p{2.9cm}p{7.5cm}}
  \toprule 编号 & 功能需求 & 需求描述与验收准则 \\\midrule\endhead
  FR-01 & 活动浏览与检索 & 访客和登录用户可查看已发布活动列表；支持关键词、分类、状态、排序和分页；详情页展示标题、时间、地点、主办方、正文、附件或来源等信息。 \\
  FR-02 & 身份认证与角色授权 & 系统提供验证码登录、邮箱验证码注册和 JWT 会话；未登录访问受限资源返回 401，角色不匹配返回 403；前端路由和后端接口均执行权限控制。 \\
  FR-03 & 活动发布与审核 & 发布者可创建和编辑草稿，提交后进入 \texttt{pending\_review}；管理员可通过或驳回，状态至少包括 \texttt{draft}、\texttt{pending\_review}、\texttt{published}、\texttt{rejected}；审核意见可回显。 \\
  FR-04 & 个人参与管理 & 登录用户可报名、取消报名、收藏、取消收藏、加入/移出个人日历，并导出 ICS；同一用户对同一活动的报名、收藏和日历关系保持幂等。 \\
  FR-05 & 发布者申请管理 & 普通用户可提交发布者申请；管理员查看申请材料并审批；审批通过后用户角色变更为 \texttt{publisher}。 \\
  FR-06 & 知识关联与订阅通知 & 活动可关联知识节点或关键词；用户可订阅节点/关键词；系统在匹配活动产生时生成站内通知，并支持通知查看和已读标记。 \\
  FR-07 & 数据源与内容质量管理 & 管理员可维护数据源，查看抓取日志；活动可记录来源、质量分和疑似重复标识；外部 URL 抓取需校验协议、域名和内网地址，降低 SSRF 风险。 \\
  FR-08 & 审计与安全治理 & 管理员可查看关键操作审计记录；密码只保存哈希，Token 不进入 URL，敏感日志脱敏；管理类导出和高权限接口仅限管理员访问。 \\
  \bottomrule
  \end{longtable}

  \subsection{功能优先级}
  \noindent\begin{tabularx}{\textwidth}{p{2.2cm}Y Y}
  \toprule 优先级 & 功能范围 & 说明 \\\midrule
  Must & FR-01、FR-02、FR-03、FR-04、FR-05、FR-08 & 支撑项目可演示和核心业务闭环，是期末提交必须覆盖的范围。 \\
  Should & FR-06、FR-07、统一异常处理、接口契约、基础测试证据 & 强化平台差异化、可信治理和工程质量，应尽量完成并提供证据。 \\
  Could & 外部搜索、AI 增强推荐、更多统计看板、更完整的无障碍审计 & 对体验和扩展有帮助，但依赖外部服务或额外工期，可作为演化方向。 \\
  Won't 本阶段 & 移动原生 App、支付、校级统一身份认证、正式生产监控平台 & 超出课程项目当前边界，不作为本次验收必需功能。 \\
  \bottomrule
  \end{tabularx}

  \subsection{功能需求用例图演示}

  详细图例请见“系统建模报告”。

  

  \section{非功能需求}
  \subsection{性能需求}
  \begin{enumerate}
  \item \textbf{API 响应时间}：读接口（活动列表、详情、搜索）在本地/CI 环境下 P95 延迟不超过 500ms；写接口（登录、注册、报名、审核）P95 延迟不超过 1s（含邮件验证码发送）。
  \item \textbf{前端加载时间}：首页首屏加载时间在模拟校园网环境下不超过 2s；活动列表页采用分页（默认每页 20 条，最大不超过 100 条）和骨架屏加载态，避免用户等待时无反馈。
  \item \textbf{并发能力}：系统应支持至少 500 用户同时在线（预发布环境参考值），核心读接口在 50 并发请求下错误率不超过 1\%。
  \item \textbf{分页与数据量}：所有列表接口必须分页返回，单页数量由服务端限制，响应中包含 \texttt{items}、\texttt{page}、\texttt{per\_page}、\texttt{total} 等分页元信息。
  \item \textbf{超时与降级}：搜索接口超时设置为 10s，AI 增强和外部抓取等耗时能力超时设置为 30s；超时或服务不可达时以降级结果或明确错误提示返回，不得阻塞活动浏览、登录、报名和审核等核心路径。
  \end{enumerate}

  \subsection{易用性需求}
  \begin{enumerate}
  \item 访客应能在不登录的情况下完成活动浏览、筛选和详情查看；登录只在报名、收藏、日历、发布和管理等受限操作前触发。
  \item 主要页面在桌面和窄屏（支持最小视口宽度 360px）下保持可读，不应出现横向溢出、按钮文字遮挡或关键信息重叠。
  \item 对加载中、无数据、请求失败、未登录、越权和提交成功等状态提供明确反馈。
  \item 活动详情页应突出时间、地点、主办方、报名入口和状态，降低用户判断成本。
  \item 表单校验应靠近输入位置提示原因，提交失败时保留用户已填写内容。
  \end{enumerate}

  \subsection{可维护性需求}
  \begin{enumerate}
  \item 前端按 \texttt{views}、\texttt{components}、\texttt{api}、\texttt{stores}、\texttt{composables} 分层，避免页面直接散落接口细节。
  \item 后端采用应用工厂、蓝图、服务层和模型层结构，领域规则尽量沉淀在服务层，便于测试和复用。
  \item 需求编号、接口、页面、测试用例和缺陷记录应保持可追踪，例如 \texttt{FR-03} 对应活动发布页、提交审核接口、审核接口和发布审核测试用例。
  \item 后端服务层函数和前端业务工具函数应保持不低于 70\% 的测试覆盖率；核心服务模块（审核、权限、质量）覆盖率目标不低于 85\%。
  \item 配置项通过 \texttt{.env.example} 和环境变量管理，密钥、数据库连接、外部服务地址不硬编码进入源码。
  \item 需求变更应经过影响分析，明确受影响的页面、接口、数据结构、测试和部署配置；紧急修复发布后也应补充回归用例。
  \end{enumerate}

  \section{数据需求}
  \subsection{数据类型与结构}
  系统涉及的核心数据类型及其关系如下：
  \begin{enumerate}
  \item \textbf{用户数据：} 用户账号（用户名、邮箱、密码哈希）、角色（\texttt{viewer} / \texttt{publisher} / \texttt{admin}）、个人资料、发布者申请记录。
  \item \textbf{活动数据：} 活动基本信息（标题、摘要、正文、时间、地点、主办方）、状态（\texttt{draft} / \texttt{pending\_review} / \texttt{published} / \texttt{rejected}）、来源、质量分和审核意见。
  \item \textbf{关系数据：} 用户与活动之间的报名关系、收藏关系、日历事件关系，均由 \texttt{(user\_id, activity\_id)} 唯一约束保证幂等。
  \item \textbf{知识数据：} 知识节点（名称、类型、别名）、活动与节点的关联关系、用户订阅记录。
  \item \textbf{审计与日志数据：} 管理操作审计日志（操作者、动作、目标类型、目标 ID、时间）、数据源抓取日志、搜索日志。
  \item \textbf{通知数据：} 站内通知（接收人、内容、关联活动、已读状态、生成时间）。
  \end{enumerate}

  \subsection{数据存储}
  \begin{enumerate}
  \item 系统采用关系型数据库（PostgreSQL）作为主存储，开发环境可使用 SQLite 替代，两者通过 SQLAlchemy 抽象层切换。
  \item 数据库表结构通过 SQLAlchemy 模型定义，当前环境使用自动建表，生产环境须使用受控迁移工具（如 Alembic）。
  \item 缓存与队列数据（会话、限流、抓取任务）存储在 Redis 中，可独立部署。
  \item 容器化部署时，PostgreSQL 和 Redis 数据挂载为 Docker 卷，确保容器重启后数据不丢失。
  \end{enumerate}

  \subsection{数据量估算}
  \begin{enumerate}
  \item \textbf{用户规模：} 以中山大学在校生约 5 万人为参考，预计活跃用户约 5000--20000 人。
  \item \textbf{活动规模：} 每学期约 200--500 场活动，全年约 500--2000 条活动记录。
  \item \textbf{关系数据：} 报名、收藏、日历事件等关系数据预计约 5--20 万条。
  \item \textbf{日志规模：} 审计日志和搜索日志随使用量线性增长，单条日志不超过 1KB，年增长预计不超过 200MB。
  \end{enumerate}
  上述数据量在 PostgreSQL 和 SQLite 的承载范围内，无需专门的分布式存储方案。

  \subsection{数据安全与备份}
  \begin{enumerate}
  \item \textbf{敏感数据保护：} 密码使用哈希（如 Werkzeug 的 \texttt{generate\_password\_hash}）存储，不保存明文；JWT 不进入 URL 参数；日志中对用户邮箱、验证码等敏感信息脱敏处理。
  \item \textbf{唯一约束与幂等：} 报名、收藏、日历关系通过数据库唯一约束提供最终一致性和幂等保证；业务层在写入前同时进行校验。
  \item \textbf{备份策略：} 开发环境数据量小，可手动备份 SQLite 文件；部署环境建议对 PostgreSQL 执行每日全量备份，定期验证恢复流程。
  \item \textbf{数据隔离：} 不同环境（开发/测试/预发布/生产）使用独立数据库，互不干扰。
  \end{enumerate}

    \section{可行性分析}
  \subsection{技术可行性}

  系统采用的技术栈均为当前业界成熟方案：
  \begin{enumerate}
  \item \textbf{前端：} Vue 3 + TypeScript + Vite + Pinia 是主流前端开发组合，Vue 3 自 2020 年正式发布以来已拥有完整生态和社区支持；TypeScript 的静态类型系统可在编译阶段捕获接口字段不匹配等问题；Vite 提供秒级热更新和高效构建；Pinia 作为 Vue 官方推荐状态管理库，API 简洁且 TypeScript 支持完善。团队通过课程实践已掌握上述技术。
  \item \textbf{后端：} Flask 作为轻量级 Python Web 框架，搭配 SQLAlchemy ORM 和 JWT 扩展，能快速实现 REST API；蓝图（Blueprint）机制支持模块化路由组织，服务层模式便于业务逻辑测试。Flask 生态成熟，相关文档和社区资源丰富。
  \item \textbf{容器化与部署：} Docker + Compose 是容器编排的标准方案，团队已编写完整的 Dockerfile 和编排配置，可在 Linux 服务器上快速重建运行环境。
  \item \textbf{持续集成：} GitHub Actions 作为 CI 平台，已配置前端构建、单元测试和端到端测试工作流，确保每次提交的可验证性。
  \item \textbf{外部依赖降级：} 邮件验证码、外部搜索和 AI 服务均设计为可选能力，当服务不可达时系统以降级模式运行，不阻塞核心业务路径。
  \end{enumerate}

  综上，系统所依赖的各项技术均已成熟且在课程项目中得到验证，不存在重大技术风险。

  \subsection{设备可行性}

  \begin{enumerate}
  \item \textbf{开发设备：} 项目在普通 Windows 笔记本上完成全部开发和测试，无需专用开发硬件。
  \item \textbf{运行环境：} 开发阶段使用 SQLite 数据库和 Mock 服务，无额外基础设施依赖；部署阶段需要一台支持 Docker 的 Linux 服务器或云主机，用于运行 PostgreSQL、Redis 和 Celery 等容器化服务。
  \item \textbf{CI 环境：} GitHub Actions 提供免费的 Ubuntu 运行环境，已配置浏览器（Chromium）和 Node.js 22 环境，满足前端自动化测试需求。
  \end{enumerate}

  团队已具备全部所需的开发和运行设备，无需新增采购。

  \subsection{进度可行性}

  \begin{enumerate}
  \item \textbf{开发周期：} 课程项目周期约为 12 周，前端与后端采用并行开发模式，通过预定义的接口契约和 Mock 服务解耦联调依赖，降低串行等待时间。
  \item \textbf{优先级管理：} 采用 MoSCoW 方法划分需求优先级，Must 功能（FR-01 至 FR-05、FR-08）构成核心业务闭环，优先交付；Should 和 Could 功能按剩余时间弹性安排。
  \item \textbf{质量保障：} CI 流水线自动执行构建、类型检查和测试，减少手工回归成本；端到端测试覆盖关键用户旅程，避免核心功能在迭代中退化。
  \item \textbf{主要风险与应对：} 前后端接口字段对齐需通过契约矩阵和联调会议保证；外部服务（邮件、AI）集成若超时可推迟至后续迭代，不影响核心演示。
  \end{enumerate}

  在 MoSCoW 优先级管理和并行开发策略保障下，团队能够在课程周期内交付满足验收标准的产品。

  \end{document}
